\chapter{Introducción} % REVISADO 1
\thispagestyle{plain} 

% \cite{} 
 
\clearpage 

Al propagarse ondas lo suficientemente fuertes en un medio no lineal éstas pueden interactuar, generando así nuevas ondas con una longitud de onda distinta, y mediante una sucesión de interacciones ir poblando un amplio rango de escalas espaciales. Al estudio estadístico de las propiedades de un ensamble de ondas no lineales débilmente acopladas fuera del equilibrio se le conoce como turbulencia de ondas. \cite{falconLaboratoryExperimentsWave2010} Éste concepto surgió originalmente a manos de Peierls (1929) para el estudio de fonones en cristales anharmónicos \cite{peierls1929kinetischen}.  

Posteriormente, en los años sesenta, esta idea resurgió independientemente en el área de la oceanografía, en primer lugar por Phillips (1960),  mostrando que dos (o tres) trenes de ondas de gravedad en la superficie de un fluido pueden interactuar para generar una tercera (o cuarta) onda fantasma, cuya amplitud crece en el tiempo, mostrando que sería posible una transferencia constante de energía entre escalas a través de interacciones resonantes entre componentes discretas del espectro \cite{phillipsDynamicsUnsteadyGravity1960, phillipsDynamicsUnsteadyGravity1961a}. A partir de esto, Longuet-Higgins (1962) mostraron que sería factible observar experimentalmente este tipo de resonancias \cite{longuet-higginsResonantInteractionsTwo1962}, siendo verificado por Longuet-Higguins (1966) \cite{longuet-higginsExperimentThirdorderResonant1966} e independientemente por McGoldrick et al. (1966) \cite{mcgoldrickMeasurementsThirdorderResonant1966}. Durante esta misma época se mostró que estas interacciones puede ocurrir también en el rango donde dominan los efectos de la tensión superficial como fuerza restitutiva por sobre la gravedad \cite{longuet-higginsGenerationCapillaryWaves1963, mcgoldrickResonantInteractionsCapillarygravity1965}.  % si se consideran los efectos de la tensión superficial, relevantes a escalas chicas, también  % fue retomada independientemente en el área de la oceanografía, en primer 

El desarrollo principal de la teoría de turbulencia de ondas vino por los trabajos de Hasselmann (1962) \cite{hasselmannNonlinearEnergyTransfer1962, hasselmannNonlinearEnergyTransfer1963} con el desarrollo de una ecuación cinética para la evolución temporal de los espectros de acción de onda $n(\vec k, t)$, tales que el espectro de energía para las ondas puede escribirse como $E(\vec k, t)=\omega(\vec k)n(\vec k,t)$, y a partir de los cuales . 

















%La turbulencia de ondas estudia la dinámica estadística de sistemas de ondas no lineales débilmente acopladas. \cite{falkovichTurbulence2006} Este concepto es completamente general, siendo aplicable a una amplia variedad de sistemas físicos, desde fonones en cristales anharmónicos, donde surgió originalmente, hasta ondas planetarias de Rossby, plasmas, condensados de Bose-Einstein, entre muchas otras. \cite{kartashovaTheoryDiscreteMesoscopic}% F
%
%En cuanto a la turbulencia de ondas gravito-capilares en la superficie libre de un fluido en contacto con la atmósfera, surgió originalmente como un intento de describir el estado del  resulta fundamental a la hora de entender los mecanismos de transferencia de energía del viento % TR 
% 
%
%
%
%
%La Tesis se organiza de la siguiente manera: el Capítulo 1 es esta introducción; el Capítulo 2 presenta las bases teóricas fundamentales en el desarrollo de Turbulencia de Ondas, haciendo énfasis en el caso particular de ondas gravito-capilares; el Capítulo 3 está dedicado al modelado, diseño, construcción y caracterización de un sensor capacitivo para la superficie libre de un fluido de alta tasa de adquisición temporal y resolución submilimétrica; el Capítulo 4 presenta un estudio experimental de turbulencia de ondas, logrado gracias a la utilización del sensor capacitivo, exponiendo los principales resultados sobre las cascadas de energía e interacciones resonantes en el flujo; finalmente el Capítulo 5 contiene las conclusiones del trabajo realizado y perspectivas a futuro. 